%%% File:        b01_Introduction.tex
%%% Author:      Joaquín Ataz-López et al
%%% Begun:       March 2020
%%% Concluded:   March 2020
%%% Contents:    First chapter of the introduction to ConTeXt: a general
%%%              overview of the system. The contents were partly
%%%              by the presentation of LaTeX
%%%              by Kopka and Daly in Chapter 1 of their Guide to
%%%              LaTeX
%%%
%%% Edited with: Emacs + AUCTeX - And at times with vim + context-plugin
%%%

\environment ../introCTX_env.tex

\startcomponent b01_Introduction.tex

\startchapter 
  [
    title=\ConTeXt:~概览,
    reference=cap:panorama
  ]

\TocChap

\startsection 
  [title=那么，什么是 \ConTeXt？]

\ConTeXt\ 是一个{\em 排版系统}，或者换句话说：一套广泛的工具集，旨在让用户能够绝对且完全地控制特定电子文档（用于纸上打印或屏幕显示）的外观和呈现。本章将解释这意味着什么。但首先，让我们强调一下 \ConTeXt\ 的一些特点。

\startitemize

  \item \ConTeXt\ 有三个{\em 版本}：Mark~II、Mark~IV 和 LMTX\null。\ConTeXt\ Mark~II 已停止开发；即，它不再处于开发中（并且已经很多年了），并且不太可能有进一步的更改或添加新功能。只有在需要更正某些错误时才会出现新版本。
    \ConTeXt\ Mark~IV 更新得多，但现在也或多或少处于停止开发状态。而 \ConTeXt\ LMTX 则仍在开发中，并且不时会出现引入一些改进或额外实用程序的新版本。尽管 LMTX 仍在开发中，但它已经相当成熟，新版本引入的更改不太可能改变现有功能的工作方式；一些更改提供了新功能，而其他更改则非常微妙，仅影响系统的底层操作。对于普通用户来说，这些更改通常是完全透明的；就好像它们没有发生一样。尽管所有三个{\em 版本}有很多共同点，但它们之间也存在一些不兼容的特性。因此，本入门只关注 \ConTeXt\ LMTX，从现在开始，不带 \quotation{Mark} 的 \quotation{\ConTeXt} 指的就是 \ConTeXt\ LMTX。

  \item \ConTeXt\ 是{\em 自由}软件（或免费软件，但不仅仅是{\em 免费}的意思）。程序本身（即构成 \ConTeXt\ 的计算机工具集）在 {\em GNU 通用公共许可证}下分发。文档则在允许自由复制和分发的 \quotation{{\em 知识共享}}许可证下提供。

  \item \ConTeXt\ 既不是文字处理器程序，也不是文本编辑程序，而是一套旨在{\em 转换}我们之前用喜欢的文本编辑器写好的文本的工具。因此，当我们使用 \ConTeXt\ 工作时：

  \startitemize

    \item 我们首先使用任何 \quotation{纯文本} 文本编辑器编写一个或多个文本文件。具体来说，这不包括文字处理器。

    \item 在这些文件中，除了构成文档内容的文本之外，还有一系列指令，告诉 \ConTeXt\ 从原始文本文件生成的最终文档必须具有的外观。完整的 \ConTeXt\ 指令集实际上是一种{\em 语言}；由于这种语言允许人们{\em 编程}实现文本的排版转换，我们可以说 \ConTeXt\ 是一种{\em 排版编程语言}。

    \item 一旦我们编写了源文件，它们将由一个程序（也称为 \MyKey{context}\footnote{\ConTeXt\ 既是一种语言也是一个程序（此外还有其他组件）。这一特性在当前文本中带来了一个问题：我们需要区分这两个方面。为此，我采用了一个排版约定：当专门指代语言，或同时指代语言和程序时，使用带有标志的写法 \quotation{ConTeXt}（\ConTeXt）；而当专门指代程序时，则使用小写等宽字体写成 \MyKey{context}。对于该语言的命令和示例，同样使用这种字体。}）处理，该程序将根据它们生成一个 PDF 文件，准备好发送到打印机、印刷厂或在屏幕上查看。

  \stopitemize

  \item 因此，在 \ConTeXt\ 中，我们必须区分我们正在编写的文档和 \ConTeXt\ 生成的文档。为避免任何疑问，在本入门中，我将包含格式化指令的文本文档称为{\em 源文件}，而将 \ConTeXt\ 根据源文件生成的 PDF 文档称为{\em 最终文档}。

\stopitemize

以上基本要点将在下面进一步展开。

\stopsection

\startsection [title=排版文本]

编写文档（书籍、文章、章节、传单、打印输出、论文 \unknown）和对其进行排版是两种截然不同的活动。编写文档在很大程度上等同于起草；这由决定其内容和结构的作者完成。作者直接创建的文档，正如他或她所写的那样，被称为{\em 手稿}。就其本质而言，只有作者或被允许阅读的人才能接触到手稿。要超越这个私密群体进行传播，手稿需要被{\em 出版}。今天，出版某物——在其 \quotation{向公众开放} 的词源意义上——就像把它放在互联网上一样简单，任何找到它并想阅读的人都可以访问。但直到相对最近，出版还是一个成本高昂的过程，依赖于某些专业化的专业人士，只有那些因其内容或作者而被认为特别有趣的手稿才能进入这个过程。即使在今天，我们也倾向于将{\em 出版}一词保留用于这种{\em 专业出版}，其中手稿的外观经历一系列旨在提高文档{\em 可读性}的转换。这一系列转换就是我们所说的{\em 排版}。

排版的目标通常是——撇开旨在吸引读者注意力的广告类文本不谈——制作具有最大{\em 可读性}的文档，即印刷文本的质量能够吸引或促进阅读，并确保读者感到舒适。许多因素促成了这一点；有些当然与文档的{\em 内容}有关（质量、清晰度、组织\unknown），但其他的则取决于所用字体的类型和尺寸、文档中空白的使用、段落之间的视觉分隔等。此外，还有其他类型的资源，不完全是图形或视觉类的，例如文档中是否存在对读者的特定辅助，如页眉和页脚、索引、词汇表、粗体的使用、边注等。了解并正确处理排版者可用的所有资源的知识可以称为 \quotation{排版艺术}或 \quotation{印刷艺术}。

历史上，直到计算机出现之前，作者和排版员的职责是相当分开的。作者用手写或称为打字机的 19 世纪机器写作，后者的排版资源甚至比手写还要有限；然后作者将原稿交给出版商或印刷商，他们将其转换以获得印刷文档。

今天，计算机科学使作者更容易决定排版的每一个细节。然而，这并不能消除这样一个事实：一个好作者所需的品质与一个好排版员所需的品质并不相同。根据所处理文档的类型，作者需要理解所写主题、清晰的阐述、结构良好的思维以创建组织有序的文本、创造力、节奏感等。但排版员必须结合对所掌握的概念和图形资源的良好了解，以及足够好的品味来和谐地使用它们。

使用一个好的文字处理程序\footnote{根据一个相当古老的惯例，我们区分{\em 文本编辑器}和{\em 文字处理器}。早期的文本编辑程序处理未格式化的文本文件，而另一种则处理格式化文本的二进制文件。}有可能获得一个排版上相当不错的文档。但一般来说，文字处理器并非为排版而设计，其结果虽然可能正确，但无法与专门为控制文档排版而设计的其他工具所获得的结果相媲美。事实上，文字处理器是打字机的演变方式，并且由于这些工具掩盖了文本创作与其排版之间的区别，它们的使用往往会产生结构不良且排版不充分的文本。相反，像 \ConTeXt\ 这样的工具是从印刷机演变而来的；它们提供了更多的排版可能性，最重要的是，不可能在学习使用它们的同时不沿途获得许多与排版相关的概念。这就是与文字处理器的区别，有些人可以使用文字处理器多年却不学一点排版知识。

\stopsection


\startsection [title=标记语言]

在计算机出现之前，如前所述，作者用手写或打字机准备手稿，然后交给出版商或印刷商，后者负责将手稿转换为最终的印刷文本。尽管作者在转换过程中的参与相对较少，但他或她确实通过一些方式进行干预，例如，指出手稿的某些行是其不同部分（章节、小节\unknown）的标题，或者指示某些内容应以某种方式在排版上加以强调。这些指示由作者直接在手稿上做出，有时是明确地，有时是通过随时间不断发展的某些惯例。
% TODO: 本段其余部分有点令人困惑。我想知道是否从西班牙语翻译丢失了一些东西。
例如，章节总是在新页开始，通过在标题前插入几个空行、在其下划线、用大写字母书写，或者将要强调的文本放在两个下划线之间、增加段落的缩进等。

简而言之，作者{\em 标记}文本以提供有关应如何排版的指示。然后，编辑会在文本上手写其他指示给印刷商，例如，要使用的字体及其大小。

今天，在计算机化的世界中，我们可以通过所谓的{\em 标记语言}继续这样做以生成电子文档。这类语言使用一系列{\em 标记}或指示，处理包含它们的文件的程序知道如何解释这些标记。如今最著名的标记语言可能是 HTML，因为今天大多数网页都基于它。一个 HTML 页面包含网页的文本，以及一系列告诉加载页面的浏览器程序应如何显示它的标记。Web 浏览器理解的 HTML 标记，连同关于如何以及在何处使用它们的指令，被称为 \quotation{HTML 语言}，这是一种{\em 标记语言}。但除了 HTML，还有许多其他标记语言；事实上它们正在蓬勃发展，因此 XML，作为{\em 首屈一指}的标记语言，如今随处可见，并且几乎用于所有事情：用于数据库设计、创建特定语言、传输结构化数据、应用程序配置文件等。还有用于图形设计（SVG、TikZ 或 MetaPost）、数学公式（MathML）、音乐（Lilypond 和 MusicXML）、金融、地理信息学等的标记语言。当然，也有旨在进行文本排版转换的标记语言，其中 \TeX\ 及其衍生品最为突出。

关于指示文本外观的{\em 排版}标记，有两种类型可以参考：{\em 纯排版标记}和{\em 概念标记}，或者如果你愿意，也可以说是{\em 逻辑标记}。纯排版标记仅限于精确指示应使用哪种排版资源来显示特定文本；例如，当我们指示某些文本应为粗体或斜体时。另一方面，概念标记指示文档中某段文本在整体中{\em 承担}什么功能，例如当我们指示某物是标题、副标题或引用时。通常，倾向于使用第二种标记的文档更加一致且更容易排版，因为它们再次指出了创作与排版之间的区别：作者指示某一行是标题，或者某个片段是警告或引用；而排版员决定如何在排版上强调所有标题、警告或引用。因此，一方面保证了一致性，因为所有{\em 承担}相同功能的片段看起来都一样，另一方面节省了时间，因为每种片段的格式只需指定一次。

\stopsection


\startsection [title=\TeX\ 及其衍生品]

\TeX\ 是在 1970 年代末由 {\sc Donald E.~Knuth} 开发的。他是斯坦福大学理论计算机编程教授（现为荣誉退休教授），开发该程序的目的有两个：制作自己的出版物，以及作为系统开发和注释的示例。\TeX\ 的开发还伴随着 \MetaFont 语言，这是一种用于设计排版字体的附加编程语言。{\sc Knuth} 使用 \MetaFont 设计了名为 {\em Computer Modern} 的字体，该字体除了包含常规字符外，还提供了一套完整的数学符号\footnote{在排版中，字形是一个字符、多个字符或部分字符的图形表示，相当于今天的字母类型（刻有字母或活字的 bits）}。结合这些工具和其他实用程序，形成了 \TeX\ 排版系统。由于其强大功能、高质量输出、灵活性和广泛可能性，\TeX\ 被认为是文本排版领域的最佳计算机化系统之一。虽然最初设计用于包含大量数学公式的文本，但人们很快发现该系统适用于所有类型的文本。

\reference[ref:boxes]{}在内部，\TeX\ 的工作方式与以前印刷厂的排版工相同。对 \TeX\ 来说，一切都是{\em 盒子}：字母包含在盒子中，空格也是盒子，几个字母（包含几个字母的盒子）形成一个包含单词的新盒子，几个单词以及它们之间的空格形成一个包含一行的盒子，几行变成一个包含段落的盒子\unknown\ 依此类推。此外，所有这些都具有非凡的测量精度。考虑到 \TeX\ 处理的最小单位比通常用于测量字符和行（这是大多数文字处理程序处理的最小单位）的排版点还要小 65,536 倍。这意味着 \TeX\ 处理的最小单位大约为 0.000005356 毫米。{\sc Knuth} 认为这对于所有实际目的来说都足够精确。

% 我从 Pablo Rodríguez 的 "Aprender ConTeXt" 中复制粘贴了带重音的 epsilon，但由于某种原因它无法处理。
% 因此我使用 \definecharacter 来创建一个带重音的 epsilon

\definecharacter etilde {\buildtextaccent ´ {\lower.2ex\hbox{\epsilon}}}

\TeX\ 这个名字来源于希腊单词 \tau\etilde\chi\nu\eta 的词根，用大写字母书写（{\switchtobodyfont[mordern]\tfx ΤÉΧΝΗ}）。因此，单词 \TeX\ 的最后一个字母不是拉丁字母 \quote{X}，而是希腊字母 \quote{\chi}，其发音显然像苏格兰语 {\em loch} 中的 \quote{ch}。所以 \TeX\ 应该发音为 {\em Tech}。另一方面，这个希腊词同时意味着 \quotation{艺术} 和 \quotation{技术}，这就是 {\sc Knuth} 选择它来命名他的系统的原因。他写道，这个名字的目的是 \quotation{提醒你 \TeX\ 主要关注高质量的技术手稿。它强调的是艺术和技术，正如其希腊词根所蕴含的意义}。

% 不要让这一行单独留在页面底部！
\need 3\baselineskip 
按照 {\sc Knuth} 建立的惯例，\TeX\ 应该写成：

\startitemize

  \item 在像本文这样的排版格式化文本中，使用我一直以来使用的标志：三个大写字母，中间的 \quote{E} 稍微向下偏移，以使 \quote{T} 和 \quote{X} 对齐更紧密；换句话说：\quotation{\TeX}。

  \startSmallPrint

    为了便于书写这个标志，{\sc Knuth} 在 \TeX\ 中包含了一个用于在最终文档中书写该标志的指令：\PlaceMacro{TeX}\tex{TeX}。

  \stopSmallPrint

  \item 在未格式化的文本中（例如电子邮件或文本文件），使用大写的 \quote{T} 和 \quote{X}，以及小写的中间 \quote{e}；因此：\quotation{TeX}。

\stopitemize

这个惯例继续用于所有包含其专有名称的 \TeX\ 衍生品中，\ConTeXt 就是这种情况。在未格式化的文本中书写时，我们需要写成 \quotation{ConTeXt}。


\startsubsection [reference=sec:engines,title=\TeX\ 引擎]

\TeX\ 程序是自由的 {\em 自由} 软件：其源代码对公众开放，任何人都可以随意使用或修改它，唯一的条件是，如果进行了修改，则结果不能称为 \quotation{\TeX}。这就是为什么随着时间的推移，出现了该程序的某些改编版本，对其进行了不同的改进，这些通常被称为 {\em \TeX\ 引擎}。除了原始的 \TeX\ 程序，按出现时间顺序，主要的引擎有 \pdfTeX、\eTeX、\XeTeX、\LuaTeX\ 和 \LuaMetaTeX\null。它们每个都应包含前一个的改进。另一方面，这些改进，在 \LuaTeX\ 出现之前，并没有影响语言本身，而只影响输入文件、输出文件、源处理以及宏的底层操作。

\startSmallPrint

  使用哪个 \TeX\ 引擎是 \TeX\ 社区内一个争论很多的问题。我在这里不展开这个问题，因为 \ConTeXt\ LMTX 只与 \LuaMetaTeX\null 一起工作。实际上，在 \ConTeXt\ 世界中，关于 \TeX\ {\em 引擎} 的讨论变成了关于使用 Mark~II（与 \pdfTeX\ 和 \XeTeX\ 一起工作）、Mark~IV（仅与 \LuaTeX\ 一起工作）还是 LMTX（仅与 \LuaMetaTeX\null 一起工作）的讨论。

\stopSmallPrint

\stopsubsection


\startsubsection [title=从 \TeX\ 衍生出的格式]

\TeX\ 的核心只理解大约 300 条非常基本的指令，称为{\em 原语}，适用于排版操作和编程功能。这些指令中的绝大多数都是非常{\em 低级}的，用计算机术语来说，这意味着它们对计算机来说比人类更容易理解，因为它们涉及非常基本的操作，比如 \quotation{将此字符向上移动 0.000725 毫米}这类。因此，{\sc Knuth} 认为 \TeX\ 应该是可扩展的，意思是应该有一种机制允许在更高级别定义指令，这些指令更容易被人类理解。这些指令在执行时被分解为其他更简单的指令，被称为{\em 宏}。例如，打印 (\tex{TeX}) 标志的 \TeX\ 指令，在执行时分解如下：

\vbox{\starttyping
T
\kern -.1667em 
\lower .5ex
\hbox {E}
\kern -.125em
X
\stoptyping}

但对于人类来说，理解和记住简单的命令 \quotation{\PlaceMacro{TeX}\type{\TeX}} 会执行打印标志所需的排版操作要容易得多。

\startSmallPrint

  什么是{\em 宏}和什么是{\em 原语}之间的区别，实际上仅从 \TeX\ 开发者的角度来看才重要。从用户的角度来看，它们是{\em 指令}，或者如果你愿意，也可以说是{\em 命令}。{\sc Knuth} 称它们为{\em 控制序列}。

\stopSmallPrint

这种通过{\em 宏}扩展语言的可能性是使 \TeX\ 成为如此强大工具的特性之一。事实上，{\sc Knuth} 自己创建了大约 600 个宏，与 300 个原语一起构成了称为 \quotation{Plain \TeX} 的系统。将 \TeX\ 本身与 Plain \TeX\ 混淆是很常见的，事实上，通常关于 \TeX\ 的几乎所有书面或口头内容，实际上都是指 Plain \TeX\null。声称是关于 \TeX\ 的书籍（包括奠基性的 \quotation{{\em The \TeX Book}}）实际上指的是 Plain \TeX；而那些认为他们直接在使用 \TeX\ 的人实际上是在使用 Plain \TeX。

Plain \TeX\ 在 \TeX\ 术语中被称为{\em 格式}，由一套广泛的宏以及关于如何使用它们的方式和方法的某些语法规则组成。除了 Plain \TeX，随着时间的推移，其他{\em 格式}也被开发出来，其中值得一提的是 \LaTeX，这是由 {\sc Leslie Lamport} 于 1985 年创建的用于 \TeX\ 的一套广泛宏，它可能是学术、技术和数学世界中使用最多的 \TeX\ 衍生品。\ConTeXt\ 作为从 \TeX\ 衍生出的格式，与 \LaTeX\ 是同一级别的。

通常，这些{\em 格式}伴随有一个程序，该程序在调用 \MyKey{tex}（或实际用于处理的引擎）处理源文件之前，将构成它们的宏加载到内存中。但是，尽管所有这些格式实际上都在运行 \TeX，但从用户的角度来看，由于每种格式都有不同的指令和不同的语法规则，我们可以将它们视为{\em 不同的语言}。它们都从 \TeX\ 中汲取灵感，但不同于 \TeX，也彼此不同。

\stopsection


\startsection [title=\ConTeXt, reference=sec:ctx]

实际上，\ConTeXt\ 最初是 \TeX\ 的一种{\em 格式}，但今天远不止于此。\ConTeXt\ 包括：

\startitemize[n]

% TODO... 3500 是正确的数字吗？
  \item 一套非常广泛的 \TeX\ 宏。如果 Plain \TeX\ 有大约 900 条指令，\ConTeXt\ 则有大约 3500 条；如果算上这些命令支持的不同选项的名称，我们谈论的是大约 4000 个单词的词汇量。词汇量如此之大是因为 \ConTeXt\ 促进学习的策略，该策略意味着包含任意数量的命令和选项同义词。

  \startSmallPrint

    其意图是，如果要实现某种效果，那么对于英语使用者称呼该效果的每种方式，都有一个命令或选项可以实现它——这应该使语言的使用更容易。例如，为了同时获得粗体和斜体字母，\ConTeXt\ 有三个指令都可以达到相同的结果：\tex{bi}、\tex{italicbold} 和 \tex{bolditalic}。

  \stopSmallPrint

  \item 同样广泛的用于 \MetaPost 的宏集，\MetaPost\ 是一种源自 \MetaFont\ 的图形编程语言，而 \MetaFont\ 又是 {\sc Knuth} 与 \TeX\ 联合开发的字体设计语言。

  \item 用 {\sc Perl}（最古老的）、{\sc Ruby}（有些也老了，有些没那么老）和 {\sc Lua}（最新的）开发的各种{\em 脚本}。

  \item 一个集成 \TeX、\MetaPost、{\sc Lua} 和 XML 的接口，允许用这些语言中的任何一种编写和处理文档，或者混合使用其中一些语言的元素。

\stopitemize

\startSmallPrint

  也许你对上面的解释不太理解？别担心！我用了很多计算机术语，提到了许多程序和语言。使用 \ConTeXt\ 并不需要了解所有这些不同的组件。重要的是，在这个学习阶段，要记住 \ConTeXt\ 集成了许多来自不同来源的工具，这些工具共同构成了一个{\em 排版系统}。

\stopSmallPrint

正是由于这后一种集成不同来源工具的特性，我们说 \ConTeXt\ 是一种旨在排版文档的 \quotation{混合技术}。我的理解是，这使 \ConTeXt\ 成为一个非常先进和强大的系统。

尽管 \ConTeXt\ 远不止是一套用于 \TeX\ 的宏集合，但它仍然基于 \TeX，这就是为什么这份我声称不过是{\em 入门}的文档侧重于这一点。

另一方面，\ConTeXt\ 比 \TeX\null 更现代。当 \TeX\ 被创建时，计算机向大众普及才刚刚开始，我们远未看到互联网和多媒体世界将会是什么样子。在这方面，\ConTeXt\ 自然地集成了一些在 \TeX\ 中一直有点格格不入的东西，例如包含外部图形、处理颜色、电子文档中的超链接、采用适合屏幕显示文档的页面尺寸等。

\stopsubsection

\startsubsection 
  [ 
    reference=sec:historyctx, 
    title=\ConTeXt\ 简史
  ]

\ConTeXt{} 大约诞生于 1991 年。它是由 {\sc Hans Hagen} 和 {\sc Ton Otten} 在一家名为 \quotation{{\em Pragma Advanced Document Engineering}} 的荷兰文档设计和处理公司创建的，通常缩写为 Pragma ADE\null。它最初是一套带有荷兰语名称的 \TeX\ 宏集合，非正式地称为 {\em Pragmatex}，面向公司中必须处理排版文档编辑的许多细节、不习惯使用标记语言或荷兰语以外的接口的非技术员工。因此，\ConTeXt{} 的第一个版本是用荷兰语编写的。其想法是创建足够数量的具有统一且一致接口的宏。大约在 1994 年，该{\em 包}足够稳定，可以用荷兰语编写用户手册，并在 1996 年，通过 {\sc Hans Hagen} 的倡议，对它的引用开始采用 \quotation{\ConTeXt{}} 这个名称。这个名字声称的意思是 \quotation{用 \TeX\ 处理文本}（使用拉丁语介词 \quote{con} 意为 \quote{与}），但同时也是一个关于英语（和荷兰语）单词 \quotation{context} 的文字游戏。因此，在这个名字背后，隐藏着一个涉及 \quotation{\TeX}、\quotation{text} 和 \quotation{context} 的三重文字游戏。

\startSmallPrint

  因此，由于这个名字基于文字游戏，\ConTeXt\ 应该发音为 \quote{context} 而不是 \quote{contecht}，因为后者会失去这个文字游戏。

\stopSmallPrint

接口大约在 2005 年开始被翻译成英语，从而产生了被称为 \ConTeXt\ Mark~II 的版本，其中 \quote{II} 的解释是，在开发者的心目中，之前的荷兰语版本是版本 \quote{I}，尽管它从未被正式这样称呼。在接口被翻译成英语后，该系统的使用开始传播到荷兰以外，接口也被翻译成其他欧洲语言，如法语、德语、意大利语和罗马尼亚语。尽管如此，\ConTeXt{} 的 \quotation{官方}文档通常基于英语版本，本文档使用的就是这个版本。

在其初始版本中，\ConTeXt\ Mark~II 与 \pdfTeX\ {\em 引擎}一起工作。但后来，随着 \XeTeX\ {\em 引擎}的出现，\ConTeXt\ Mark~II 被修改以允许使用这个新引擎，与 \pdfTeX\null 相比，它带来了许多优势。但是当几年后 \LuaTeX\ 出现时，决定在内部重新配置 \ConTeXt{} 的工作方式，以集成这个新引擎提供的所有新可能性。于是，\ConTeXt\ Mark~IV 诞生了，它于 2007 年推出，紧随 \LuaTeX\null 推出之后。很可能，影响重新配置 \ConTeXt\ 以适应 \LuaTeX\ 的决定的因素之一是，\ConTeXt{} 的三位主要开发者中的两位，{\sc Hans Hagen} 和 {\sc Taco Hoekwater}，同时也是开发 \LuaTeX\null 的主要团队的一部分。这就是为什么 \ConTeXt\ Mark~IV 和 \LuaTeX\ 同时诞生并协同发展。\ConTeXt{} 和 \LuaTeX\ 之间存在一种在 \TeX\ 的任何其他衍生品中都不存在的协同作用；我怀疑是否有任何其他衍生品能像 \ConTeXt{} 那样利用 \LuaTeX\ 的优势。

% TODO: 写一些关于 luametatex 和 LMTX 诞生的内容。

Mark~II、Mark~IV 和 LMTX 之间有许多差异，尽管其中大部分是{\em 内部}的，也就是说，它们与宏在更低级别实际工作的方式有关，因此从用户的角度来看，差异并不明显：许多宏的名称和参数保持不变。然而，确实存在一些影响接口的差异，迫使人们根据所使用的版本以不同的方式做事。这些差异相对较少，但它们确实影响到非常重要的方面，例如输入文件的编码或系统安装字体的处理。话虽如此，在 2026 年，Mark~II 和 Mark~IV 应被视为历史工具，不应用于开发新文档。

\startSmallPrint

% TODO: 这对 Wiki 的评价仍然准确吗？
  然而，如果某处有一份文档解释（或列出）Mark~II、\Conjecture Mark~IV 和 LMTX\null 之间明显的差异，那将是非常受欢迎的。例如，在 \ConTeXt\ wiki 中，对于每个 \ConTeXt\ 命令，有时会有{\em 两种语法}（通常完全相同）。我推测一种属于 Mark~II，另一种属于 Mark~IV 和/或 LMTX；基于这个假设，我还推测{\em 第一个版本}来自 Mark~II\null。但事实是 wiki 对此只字未提。

\stopSmallPrint

事实上，在语言层面上差异相对较少，这意味着在许多情况下，我们谈论的不是两个版本，而是 \ConTeXt{} 的两种 \quotation{风味}。但无论你怎么称呼它们，事实是，为一种 \ConTeXt\ 风味准备的文档通常不能用另一种风味编译；如果文档混合了两种或三种版本，它很可能无法与其中任何一种很好地编译。这意味着源文件的作者必须首先决定是为 Mark~II、Mark~IV 还是 LMTX\null 编写。（自从 LMTX 出现以来，很少有或根本没有理由为 Mark~II 或 Mark~IV 编写新文档；但是，如果想要编辑为这两种风味中的任何一种编写的现有文档，那么使用原始文档风味的命令 {\it 可能} 比重写文档以使用适合 LMTX 的命令更容易。）

\startSmallPrint

  如果我们使用不同版本的 \ConTeXt{}，一个区分 Mark~II 和 Mark~IV 文件的好技巧是使用不同的文件扩展名。因此，例如，对于我为 Mark~II 编写的任何文件，我使用 \MyKey{.mkii} 作为扩展名；对于 Mark~IV，我使用 \MyKey{.mkiv}，而对于为 LMTX\null 编写的文件，我使用 \MyKey{.tex}。诚然，\ConTeXt{} 期望所有源文件都有扩展名 \MyKey{.tex}，但只要我们在将 \ConTeXt{} 应用于文件时明确指定文件扩展名，就可以更改文件扩展名。

\stopSmallPrint

截至 2026 年，\suite- 发行版包括 Mark~IV 和 LMTX，两者都可以使用命令 \MyKey{context} 访问。Mark~II 仍然可用（主要供需要处理旧文档的人使用），但它不包含在 \suite- 发行版中。然而，著名的 \TeX{}live 发行版包含了 Mark~II，可以使用命令 \MyKey{texexec} 访问。

% TODO 我不知道这点的真实性，我认为无论如何可以省略。
%%\startSmallPrint
%%
%%  事实上，命令 \MyKey{context} 和 \MyKey{texexec} 都是带有不同选项的{\em 脚本}，它们运行 \MyKey{mtxrun}，而 \MyKey{mtxrun} 又是一个 {\sc Lua} {\em 脚本}。
%%
%%\stopSmallPrint

如今，Mark~II 和 Mark~IV 已冻结，而 LMTX 仍在开发中，这意味着只有发现需要更正的错误或缺陷时才会发布旧版本的新版本，而 LMTX 的新版本则继续定期发布；有时一个月两到三次，尽管在大多数情况下，这些 \quotation{新版本} 并没有引入语言上可感知的变化，而是仅限于在低级别上改进命令的实现方式，或更新发行版中包含的众多手册中的一些。即便如此，如果我们安装了开发版本——这是我推荐的，也是 \suite-{} 默认安装的版本——定期更新我们的版本是有意义的（参见 \in{附录}[installation_suite] 了解更新已安装 \suite- 版本的方法）。

% TODO 大部分内容已被上面的更改所包含。
%      但应审查以下内容，看看其中哪些部分应该在此呈现。
%% \startSmallPrint
%% 
%% \startsubsubsubsubject 
%%   [title=LMTX 和 Mark~IV 的其他替代实现]
%% 
%% \ConTeXt{} 的开发者自然是 restless 的，因此他们并没有停止对 Mark~IV 的 \ConTeXt{} 开发；新版本仍在测试和试验中，尽管通常这些版本与 Mark~IV 的差异很小，并且不存在 Mark~IV 和 Mark~II 之间那样的编译不兼容性。
%% 
%% 因此，出现了某些称为 Mark~VI、Mark~IX 和 Mark~XI 的 Mark~IV 小变体。其中，我只能在 \ConTeXt{} wiki 中找到关于 Mark~VI 的一个小引用，其中说与 Mark~IV 的唯一区别在于定义命令的可能性，不是像传统 \TeX\ 那样为参数分配一个数字，而是分配一个名称，就像几乎所有编程语言中通常所做的那样。
%% 
%% 我认为，比这些小变体更重要的是，\ConTeXt{} 领域中出现了一个名为 LMTX 的新版本，这个名字是 LuaMetaTeX 的首字母缩写：一个全新的 \TeX\ {\em 引擎}，是 \LuaTeX\ 的简化版本，旨在节省计算机资源；这意味着 LMTX 比 \ConTeXt\ Mark~IV 需要更少的内存和处理能力。
%% 
%% LMTX 于 2019 年春季推出，人们认为它不会对 Mark~IV 语言带来任何外部变化。对于文档作者来说，使用它时不会有任何区别；但在编译时，需要选择使用 \LuaTeX\ 还是使用 LuaMetaTeX。在关于安装 \ConTeXt 的 \in{附录}[installation_suite] 中，展示了一个为每个安装分配不同命令名称的过程（\in{节}[sec:alias]）。
%% 
%% \stopsubsubsubject
%% 
%% \stopSmallPrint

\stopsubsection

\startsubsection [title=\ConTeXt\ 与 \LaTeX\ 的比较]

鉴于从 \TeX{} 衍生出的最流行的格式是 \LaTeX{}，这两者之间的比较是不可避免的。显然，我们谈论的是不同的语言，尽管它们在某种程度上相互关联，因为它们都源自 \TeX\null。这种关系类似于例如西班牙语和法语之间的关系：有共同起源（拉丁语）的语言，这意味着它们的语法{\em 相似}，并且每种语言中的许多单词在另一种语言中都有对应的单词。但除了这种{\em 家族相似性}之外，\LaTeX\ 和 \ConTeXt\ 在理念和实现上有所不同，因为两者的初始目标在某种程度上是相反的。\LaTeX\ 声称要简化 \TeX\ 的使用，将作者与具体的排版细节隔离开来，以帮助专注于内容，将排版细节留给 \LaTeX\null。这意味着简化 \TeX\ 的使用是以牺牲 \TeX\ 的巨大灵活性为代价的，通过预定义基本格式并限制作者必须决定的排版问题的数量。与这种理念相反，\ConTeXt\ 诞生于一家致力于排版文档的公司内部。因此，它远非想要将作者与排版细节隔离开来，而是旨在赋予作者对它们的绝对和完全控制。为了实现这一点，\ConTeXt\ 提供了一个统一且一致的接口，比 \LaTeX\null 更接近 \TeX\ 的原始精神。

这种理念和创始目标的差异进而转化为实现上的差异。\LaTeX\ 倾向于尽可能简化，不需要使用 \TeX\ 的所有资源。在某种程度上，其核心相当简单。因此，当需要扩大其可能性时，有必要明确编写一个{\em 包}来实现。这种与 \LaTeX\ 相关的{\em 打包}既是优点也是缺点：优点是因为 \LaTeX\ 的巨大知名度，加上其用户的慷慨，意味着我们可能遇到的几乎任何需求都已被某人满足过，并且有一个包可以实现它；但这也是一个缺点，因为这些包之间常常不兼容，而且它们的语法并不总是统一的。这意味着使用 \LaTeX\ 需要不断搜索成千上万个已有的包来满足需求，并确保它们能一起工作。

与通过包扩展来补充的 \LaTeX\ 核心的简单性相反，\ConTeXt\ 被设计为在其内部包含所有——或几乎所有——\TeX\ 的排版可能性，因此它的概念更加一体化，但同时也更加模块化。\ConTeXt\ 核心允许我们做几乎所有事情，并且我们保证其不同实用程序之间不会有不兼容性，无需为我们需要的功能研究扩展，语言的语法也不会仅仅因为我们需要一个特定的实用程序而改变。

诚然，\ConTeXt\ 有所谓的扩展{\em 模块}，有些人可能认为它们发挥着类似于 \LaTeX\ 包的功能，但实际上它们的工作方式不同：\ConTeXt\ 模块专门设计用于包含额外的实用程序，这些实用程序要么因为仍处于实验阶段而尚未并入核心，要么允许访问由 \ConTeXt\ 开发团队之外的某人所编写的扩展。

我不认为这两种{\em 理念}中哪一种更可取。问题更多地取决于用户的概况和需求。如果用户不想处理排版问题，只想制作非常高质量的标准化文档，那么选择像 \LaTeX\ 这样的系统可能更可取；另一方面，喜欢尝试、或需要控制文档每一个细节的用户，或者必须为文档设计特殊布局的用户，使用像 \ConTeXt\ 这样的系统可能更好，作者可以掌握所有控制权；当然，也有不知道如何正确使用这种控制的风险。

\stopsubsection


\startsubsection 
  [title=理解使用 \ConTeXt\ 的工作流程]

当我们使用 \ConTeXt\ 时，我们总是从编写一个文本文件（我们称之为{\em 源文件}）开始，在其中，除了最终文档的实际内容外，我们还将包含指示我们希望如何精确格式化文档的指令（用 \ConTeXt\ 的话说）：我们希望页面和段落具有的总体外观、我们希望应用于某些段落的边距、我们想要显示的字体、我们希望以不同字体显示的片段等。一旦我们编写了源文件，我们从终端（或者可能是我们的文本编辑器）应用 \MyKey{context} 程序，它将处理该文件，并从中生成一个不同的文件，其中文档的内容将根据源文件中为此目的包含的指令进行格式化。这个新文件可以发送给（商业）打印机、在屏幕上显示、放在互联网上或在联系人、朋友、客户、老师、学生\unknown\ 之间分发，或者换句话说，分发给任何我们为其编写文档的人。

这意味着在使用 \ConTeXt\ 时，作者处理的是一个视觉外观与最终文档毫无关系的文件：作者直接处理的文件是一个没有进行排版格式化的文本文件。因此，\ConTeXt\ 的工作方式与所谓的{\em 文字处理器}程序不同，后者在我们编写文档的同时就显示编辑后文档的最终外观。对于那些习惯于文字处理器的人来说，使用 \ConTeXt\ 的工作方式起初会感到奇怪，甚至可能需要一段时间来适应。然而，一旦习惯了，人们就会明白，这种区分工作文件和最终结果的工作方式实际上是一种优势，原因有很多，在此我无特定顺序地强调以下几点：

\startitemize[n,broad]

  \item 因为文本文件比文字处理器的二进制文件 \quote{更轻}，处理它们需要更少的计算机内存，它们更不容易损坏，并且当我们更改创建它们的程序版本时，它们不会变得不可读。它们也与任何操作系统兼容，并且可以用许多文本编辑器编辑，因此为了处理它们，计算机系统不必安装创建该文件的程序：任何其他编辑程序都可以；并且在每个计算机系统中，总会有某种文本编辑程序。

  \item 因为区分工作文档和最终文档有助于区分文档的实际内容与其外观，允许作者在创作阶段专注于内容，在排版阶段专注于外观。

  \item 因为它允许快速准确地更改文档的外观，因为这由可以轻松识别的 \ConTeXt\ 命令决定。

  \item 因为这种更改外观的便利性，另一方面，允许我们轻松地从单一内容生成两个（或更多）不同的版本：可能一个版本针对纸上打印进行了优化，另一个版本设计为在屏幕上显示，适应后者的尺寸，并可能包含在纸质文档中没有意义的超链接。

  \item 因为文字处理器中常见的排版错误（例如将斜体扩展到单词最后一个字符之后）也很容易避免。

  \item 因为当工作文件不被分发并且 \quotation{仅供我们自己查看}时，可以为我们自己后续的修订或版本加入注释和观察、评论和警告，安心地知道这些不会出现在要分发的格式化文件中。

  \item 因为同时处理整个文档所能获得的质量远高于使用必须在编写文档时做出排版决策的程序所能达到的质量。

  \item 等等。

\stopitemize

% TODO 又是 3500 这个数字。可以吗？

以上所有意味着，一方面，在使用 \ConTeXt\ 时，一旦我们掌握了窍门，我们就会更高效、更有生产力，另一方面，我们能获得的排版质量远优于使用所谓的{\em 文字处理器}所能获得的质量。尽管后者确实更容易使用，但实际上它们并没有{\em 那么}容易使用。因为虽然 \ConTeXt{} 如前所述包含 3500 条指令，但 \ConTeXt{} 的普通用户并不需要知道所有这些指令。要做通常用文字处理器做的事情，我们只需要知道允许我们指示文档结构的指令、一些关于常见排版资源的指令（如粗体或斜体），以及可能如何生成列表或脚注的指令。总共不超过 15 或 20 条指令就能让我们完成几乎所有用文字处理器能做的事情。其余的指令允许我们做不同的、通常无法用文字处理器完成或很难实现的事情。我们可以说，虽然学习使用 \ConTeXt{} 比学习使用文字处理器更困难，这是因为用 \ConTeXt{} 可以做得更多。

\stopsubsection


\startsubsection 
  [title=获取 \ConTeXt\ 的帮助]

\adaptlayout[line=+2]

在我们还不熟悉的时候，获取 \ConTeXt\ 帮助的最佳地点无疑是 \goto{wiki}[url(wiki)]，那里有很多示例，并且有很好的搜索引擎，尤其是如果你英语好的话。当然，我们也可以在互联网上找到帮助，但在这里，\ConTeXt\ 名称中的文字游戏会给我们带来麻烦，因为搜索单词 \quotation{context} 会返回数百万个结果，其中大多数与我们寻找的内容无关。要查找有关 \ConTeXt\ 的信息，你需要在 \quotation{context} 一词上添加一些内容；例如，\quotation{tex}、或 \quotation{Mark IV}、或 \quotation{Hans Hagen}（\ConTeXt\ 的创建者之一）、或 \quotation{Pragma ADE}，或类似的东西。使用 wiki 名称 \quotation{contextgarden} 搜索信息也可能有用。

当我们对 \ConTeXt\ 有了更多了解后，我们可以查阅 \suite- 中包含的众多文档中的一些，甚至在 \goto{TeX -- LaTeX Stack Exchange} [url(https://tex.stackexchange.com/)] 或 \ConTeXt\ 邮件列表（\goto{NTG-context}[url(https://mailman.ntg.nl/mailman/listinfo/ntg-context)]）上寻求帮助。后者涉及最了解 \ConTeXt\ 的人，但良好的网络礼仪规则要求，在提问之前，我们应该事先努力自己寻找答案。

\stopsubsection

\stopsubsection

\stopchapter

\stopcomponent

%%% Local Variables:
%%% mode: ConTeXt
%%% eval: (auto-fill-mode 1)
%%% TeX-master: "../introCTX_eng.tex"
%%% coding: utf-8-unix
%%% End:
%%% vim: tw=72: